\scanpage{151}
\deleted{\uline{Corollaire}}

Nous pouvons préciser le cor. 1 ainsi :

\uline{Cor 4} Sous les conditions du cor. 1, posons \inserted{pour $1\leq i\leq\unsure{r|\nu}$} \deleted{donnons nous pour tt $1\leq i\leq\unsure{r|h|\nu}$ un iso}
\[
\unsure{R|P}_i=E_i/E_{i-1},\qquad R'_i=E'_i/E'_{i-1}\qquad(E_0=0,E'_0=0)
\]
et donnons nous des isom
\[
\begin{aligned}
u_i:\ &R_i\xrightarrow{\sim}R'_i\\
u_q:\ &Q\xrightarrow{\sim}Q'
\end{aligned}
\]
avec $v$ \unsure{orthogonale|orthogonal} pour $q,q'$. Alors, \deleted{il existe} \inserted{$S$ étant affine,} \inserted{a) il existe} un isom
\[
u:(E,q)\longrightarrow(E',q')
\]
tel que $u(D)=D'$, et que $u$ induise les $u_i$ et $u_q$.

\deleted{\unsure{b)}} L'ens de ces $u$ est \inserted{(\unsure{évidemt|évident})} un torseur à dr. sous le \inserted{ss} groupe \inserted{$\underline{\mathrm O}_D(q)^+$,} des \deleted{\unsure{auto|$\to$} $\mathrm O_D(q)(S$} $\mathrm{Aut}(E,q,D)=\underline{\mathrm O}_D(q)(S)$, formé des autom qui induisent l'identité sur les $G_i$ et sur $Q$.

\deleted{Considérons d'autre} \inserted{Considérons d'autre} part l'ens \deleted{des} produit
\[
\Sigma=\prod_{1\leq i\leq r}\Sigma_i\times\Sigma_0
\]
où $\Sigma_i$ est l'ens des scindages de
\[
0\longrightarrow E_{i-1}\mathbin{-}E_{\deleted{\unsure{r}}i}\mathbin{-}R_i\mathbin{-}0
\]
et $\Sigma_0$ l'ens des scindages \uline{isotropes} de
\[
0\longrightarrow\deleted{\unsure{E}}E_r\longrightarrow E_{r\deleted{\unsure{2}}}^{\perp}\longrightarrow Q\longrightarrow0
\]
\deleted{\unsure{On}} Considérons \deleted{d'autre part} \inserted{aussi}
\[
\overline E=\bigoplus(R_i+\check R_i)\oplus Q
\]
avec la forme $\overline q$ somme des formes
% archive sheet number: 150
